\section{Convenciones de Fourier}
\label{FourierConventions}

En este trabajo se utiliza la transformada de Fourier con frecuencia medida en ciclos por unidad. Esta elección fija los factores de $2\pi$ en las exponenciales y será la misma que se usará posteriormente al interpretar las frecuencias discretas producidas por la FFT.

Inicialmente las identidades se formulan para funciones de Schwartz, lo cual permite justificar sin dificultades los cambios de variable y las integraciones que aparecen en la derivación de la retroproyección filtrada. Denotamos por $\mathcal{S}(\mathbb{R}^{d})$ el espacio de las funciones $\varphi\in C^{\infty}(\mathbb{R}^{d})$ tales que, para todo par de multiíndices $\alpha,\beta$,
\[
\sup_{x\in\mathbb{R}^{d}}\left|x^{\alpha}\partial^{\beta}\varphi(x)\right|<\infty.
\]
En particular, se usarán $\mathcal{S}(\mathbb{R})$ y $\mathcal{S}(\mathbb{R}^{2})$.

Si $h\in \mathcal{S}(\mathbb{R})$, su transformada de Fourier unidimensional se define por
\[
\widehat h(\sigma)
=
\mathcal{F}h(\sigma)
=
\int_{\mathbb{R}} h(t)e^{-2\pi i t\sigma}\,dt,
\qquad \sigma\in\mathbb{R}.
\]
La transformada inversa viene dada por
\[
h(t)
=
\mathcal{F}^{-1}\widehat h(t)
=
\int_{\mathbb{R}} \widehat h(\sigma)e^{2\pi i t\sigma}\,d\sigma.
\]
En dos variables, para $f\in \mathcal{S}(\mathbb{R}^{2})$ se define
\[
\widehat f(\xi)
=
\mathcal{F}_{2}f(\xi)
=
\int_{\mathbb{R}^{2}} f(x)e^{-2\pi i x\cdot \xi}\,dx,
\qquad \xi\in\mathbb{R}^{2},
\]
y la inversión toma la forma
\[
f(x)
=
\int_{\mathbb{R}^{2}} \widehat f(\xi)e^{2\pi i x\cdot \xi}\,d\xi.
\]
Usaremos $\sigma$ para la variable de frecuencia unidimensional, asociada a la variable del detector de las proyecciones, y $\xi$ para la variable de frecuencia en $\mathbb{R}^{2}$.

Con esta normalización, el teorema de Plancherel afirma que la transformada de Fourier se extiende de manera unitaria de $\mathcal{S}$ a $L^{2}$, es decir,
\[
\|\widehat h\|_{L^{2}(\mathbb{R})}=\|h\|_{L^{2}(\mathbb{R})},
\qquad
\|\widehat f\|_{L^{2}(\mathbb{R}^{2})}=\|f\|_{L^{2}(\mathbb{R}^{2})}.
\]
Este resultado estándar será usado para interpretar los filtros como operadores en el dominio de frecuencias; puede consultarse en Natterer \cite[Sec.~VII.1]{natterer2001} y en Kuchment \cite[Ap.~B.8]{kuchment2014}. Estas fuentes pueden emplear normalizaciones de Fourier distintas de la adoptada aquí, pero las formulaciones son equivalentes mediante un cambio de escala en la variable de frecuencia.

Es importante distinguir esta convención de la convención de frecuencia angular. Si se escribe $\eta=2\pi\sigma$, entonces las fórmulas equivalentes contienen factores adicionales de $2\pi$. En esta tesis no se usará $\eta$ como variable principal; el símbolo de los filtros se escribirá siempre como función de $\sigma$.
